The paper that precipitated this work is Goldsmith\cite{goldsmith_assessing_2016}.  
In it, he illustrates a method to model penalized function-on-scalar regression, 
permitting estimation of population-level fixed effects, subject-level random effects, and residual covariance using penalized splines.  
The working assumption for this model is that for subjects \(i=1,\dots, I\) and replications \(j=1,\dots, J_i\), \(n=\sum_i J_i\) 
functional observations recorded in one-dimensional space are available.  Each function is defined over a common domain 
\(\mathcal{T} \subset \mathbf{R}\) and is observed on a discrete grid of \(D\) evaluation points \(\{t_1, \dots, t_D\}\).  
This model, in its most basic form is written as:

\begin{equation}
\label{functional_mixed}
    \boldsymbol{y}_{ij} = \boldsymbol{\omega}_i\boldsymbol{\beta} + \boldsymbol{z}_{i}\boldsymbol{b} + \boldsymbol{\epsilon}_{ij}.
\end{equation}

\noindent
This is simply a standard linear mixed effects model: \(\boldsymbol{y_{ij}}\) denotes a vector of functional outcomes for subject 
\(\emph{i}\) and replication \(\emph{j}\), \(\boldsymbol{\omega_i}\) denotes subject specific covariates, \(\boldsymbol{\beta}\) 
denotes the fixed effect coefficient, \(\boldsymbol{z_i}\) denotes a design matrix for subjects' random effects, \(\boldsymbol{b}\) 
denotes random effect coefficients.  
The key difference from the longitudinal mixed effects model \eqref{fda:lmm} is that the fixed and random effects in 
\eqref{functional_mixed}, \(\boldsymbol{\beta}\) and \(\boldsymbol{b}\), are multivariate functions.

Following the work of Silverman\cite{ramsay_functional_2005}, Reiss\cite{reiss_fast_2010}, 
and Goldsmith\cite{goldsmith_assessing_2016}, we write \(\boldsymbol{\beta}\) and \(\boldsymbol{b}\) using spline expansions.  
They define $\boldsymbol{\beta} = \left[\Theta B_w\right]^T$ and $\boldsymbol{b} = \left[\Theta B_Z\right]^T$, where \(\Theta\) 
is the basis function matrix of dimensions \(D\times K_{\theta}\), \(K_\theta\) is the number of basis functions used, \(B_w\) 
is a matrix containing basis coefficients for \(\boldsymbol{\beta}\), and \(B_z\) is a matrix containing basis coefficients for 
\(\boldsymbol{b}\).  For a univariate model, this provides:

\begin{equation}
    \boldsymbol{y}_{ij} = \boldsymbol{\omega}_iB_W^T\Theta^T + \boldsymbol{z}_{i}B_Z^T\Theta^T + \boldsymbol{\epsilon}_{ij}.
\end{equation}

In this foundational model, Goldsmith penalizes the fixed and random effect spline coefficients with a common, 
\emph{a priori} penalty matrix\cite{goldsmith_assessing_2016}.  
Goldsmith also allows for the variances for the fixed effect coefficient to vary by covariate.  
This allows for a possibly different level of smoothness for each functional coefficient curve.  
The random effect coefficients  are penalized together (i.e. not different smoothness penalties for each subject's random effects).  
This leaves us with the following model:

\begin{equation}
    \begin{split}
        \boldsymbol{y}_{ij} = \boldsymbol{\omega}_iB_W^T\Theta^T + \boldsymbol{z}_{i}B_Z^T\Theta^T + \boldsymbol{\epsilon}_{ij}\\
        \epsilon\sim N[0, \Sigma\otimes I_n]\\
        B_{Z,i} \sim N[0, \sigma_z^2P^{-1}]\quad \text{for} \quad i = 1,\dots, I\\
        B_{W, k} \sim N[0, \sigma_{W,k}^2P^{-1}],
    \end{split}
\end{equation}

\noindent
where \(\epsilon\) denotes the stacked collection of all \(\boldsymbol{\epsilon}_{ij}\).  
Goldsmith assigned inverse-gamma priors on the variance components and an unstructured inverse-Wishart prior for the residual 
covariance matrix\cite{goldsmith_assessing_2016}.  A section of his paper is dedicated to choosing hyper-parameters for these priors; 
he uses a series of linear regressions and principal component analyses to obtain hyperparameter values a la an empirical Bayes approach.   

Goldsmith points out that the sampling performance of this model can be improved by hierarchical recentering, and, indeed, 
recentering is a common method\cite{gelfand_efficient_1995}, although Gelman refers to this as \say{decoupling}\cite{gelman_bayesian_2013}.  With this recentering/decoupling, Goldsmith presents the following model for one dimension:

\begin{equation}
    \begin{split}
        Y = ZB_Z^T\Theta^T + \epsilon\\
        \epsilon\sim N[0, \Sigma\otimes I_n]\\
        \Sigma \sim IW[v, \Psi]\\
        B_{Z,i} \sim N[B_w \boldsymbol{w}_i^T, \sigma_z^2P^{-1}]\quad \text{for} \quad i = 1,\dots, I\\
        \sigma_z^2 \sim IG[a_z, b_z]\\
        B_{W, k} \sim N[0, \sigma_{W,k}^2P^{-1}]\\
        \sigma_{W,k}^2 \sim IG[a_{W,k}, b_{W,k}]\quad \text{for} \quad k = 1,\dots, p
    \end{split}
    \label{mod1_var}
\end{equation}

Here \(Y\) is a \(n\times D\) matrix composed of functional observations stacked in rows, \(Z\) is the random effects design matrix, 
and \(W\) denotes the fixed effects design matrix.  Goldsmith then derives the conditional distributions of 
\(p\left(B_{Z,i}|\text{rest}\right)\), \(p\left(\text{vec}(B_W)|\text{rest}\right)\), \(p\left(\sigma_z^2|\text{rest}\right)\), 
\(p\left(\sigma_{W,k}^2|\text{rest}\right)\), and \(p\left(\Sigma|\text{rest}\right)\) in order to set up a Gibbs sampler. 